
\documentclass{beamer}
\usepackage[UTF8]{ctex}
\usepackage[english]{babel}
\usepackage{amsmath,amssymb}
\usepackage{hyperref}
\usepackage{graphicx}
\usepackage{xcolor}
\usetheme{Madrid}
\usecolortheme{seahorse}
\setbeamertemplate{navigation symbols}{}
\setbeamertemplate{items}[circle]
\setbeamerfont{frametitle}{size=\Large,series=\bfseries}
\definecolor{slidebg}{HTML}{FFFFFF}
\setbeamercolor{background canvas}{bg=slidebg}
\begin{document}
% ===== Slide 1: Problem and Motivation =====
\begin{frame}{Problem and Motivation}
    \textbf{Zeroth-order control for combinatorial Stackelberg congestion games.}
    \vspace{1em}
    \begin{itemize}
        \item • We propose treating combinatorial follower equilibrium solvers as nondifferentiable black boxes and optimizing leader parameters via zeroth-order randomized smoothing to robustly handle hyper-objective nonsmoothness in Stackelberg congestion games.
        \item • This avoids differentiating through large inner solvers and instead uses two-point finite-difference gradient estimators, enabling scalability when strategy sets are exponentially large or equilibria change discontinuously with parameters.
        \item • To support combinatorial linear minimization, we introduce ZDD-based exact or stratified subsampling LMOs and derive convergence and error-propagation guarantees for the Frank–Wolfe inner solver and outer ZOS optimizer.
    \end{itemize}
    \vspace{1em}
    \textit{Lead into the method overview of the inner Frank–Wolfe equilibrium solver and the outer zeroth-order optimizer.}
\end{frame}

% ===== Slide 2: Method =====
\begin{frame}{Method}
    \textbf{Inner Frank–Wolfe and outer zeroth-order loop.}
    \vspace{1em}
    \begin{itemize}
        \item • Method figure: Flow of inner Frank–Wolfe equilibrium solver and outer zeroth-order randomized-smoothing optimizer with two-point estimator.
        \item • Inner loop: projection-free Frank–Wolfe runs T iterations using a combinatorial LMO (exact shortest-path or ZDD-based sampling).
        \item • LMO: exact poly-time or ZDD DP for NP-hard cases, or subsampled candidates using uniform or length-stratified sampling.
        \item • Outer ZOS: randomized-ball smoothing, mini-batched two-point finite differences, projected updates, and parameter choices to control inner approximation error.
    \end{itemize}
    \vspace{1em}
    \textit{Next show main experimental performance versus the Diff baseline with runtime and memory comparisons.}
\end{frame}

% ===== Slide 3: Results =====
\begin{frame}[t]{Results}
    \begin{columns}[T,onlytextwidth]
        \begin{column}{0.50\textwidth}
            \textbf{Empirical speed and memory dominance over Diff baseline.}
            \vspace{0.8em}
            {\small
            \begin{itemize}
                \item • Main result: ZOS matches Diff in social cost at small FW gaps while outperforming the Diff baseline by approximately 20–1000× per-iteration speed.
                \item • ZOS uses far less memory ($\leq$1.7 GiB peak) versus Diff’s 10–194 GiB, enabling practical scaling on large ZDD instances.
            \end{itemize}
            }
            \vspace{0.8em}
            \textit{Conclude with key practical findings, trade-offs, and limitations.}
        \end{column}
        \begin{column}{0.48\textwidth}
            \centering
            \includegraphics[width=\linewidth,height=0.7\textheight,keepaspectratio]{\detokenize{../main_figure_table/table_1.png}}
        \end{column}
    \end{columns}
\end{frame}

% ===== Slide 4: Key Findings =====
\begin{frame}{Key Findings}
    \textbf{Practical guarantees and sampling trade-offs for scalable Stackelberg control.}
    \vspace{1em}
    \begin{itemize}
        \item • Treating equilibrium solvers as black boxes enables scalability and large speedups by avoiding costly backpropagation through many inner steps.
        \item • ZDD-based stratified sampling recovers accuracy with moderate sample sizes (≈1000), reduces memory, but may plateau or require large m in some regimes.
        \item • Theoretical error propagation yields Goldstein-stationarity guarantees while practical choices of T, m, B, ρ, and η balance inner variance and outer convergence.
    \end{itemize}
    \vspace{1em}
    \textit{End of presentation; invite questions about implementation details and parameter tuning.}
\end{frame}
\end{document}
